% GENERATED FILE. Edit canonical lessons, then run npm run generate:books.
% canonical-source-hash: fnv1a64:335d13b4f186eec6
% canonical-lessons: AR-C143-mimhah, AR-C143-saff, AR-C143-talib, AR-C143-wazifa, AR-C143-ratib

\chapter{Eraser, Class, Student, Job, Salary}
\label{ch:ar-eraser-class-student-job-salary}
\hlchaptermodality{143}

\begin{chapteropening}
\textbf{By the end of this chapter:} \emph{I can say an eraser, a class, a student, a job and a salary.}

Complete the last lesson of chapter 143: I can say an eraser, a class, a student, a job and a salary.
\end{chapteropening}

\section[\texorpdfstring{\ar{ممحاة} (mimḥāh)}{ممحاة (mimḥāh)}]{\ar{ممحاة} (mimḥāh) — an eraser}

\label{lesson:AR-C143-mimhah}



\begin{quote}
Before the new one: say the Arabic for a factory, then the Arabic for a square.
\end{quote}

\subsection*{You'll want to know: \ar{ممحاة}}
\textbf{\ar{ممحاة}} — \emph{mimḥāh} — \textquotedblleft{}an eraser\textquotedblright{}.

\begin{cousinweb}
\textbf{\ar{ممحاة}} (\emph{mimḥāh}) is built on the native Arabic root \emph{m-ḥ-w}, \textquotedblleft{}to efface, to erase\textquotedblright{} (Lane). \emph{Mimḥāh} is an instrument noun, \textquotedblleft{}a thing for erasing\textquotedblright{}.
\end{cousinweb}

\begin{grammarlens}[title={What you've built}]
One more word for this chapter.
\end{grammarlens}

\subsection*{Your turn}
\begin{itemize}
\raggedright
  \item \emph{Say it:} \emph{mimḥāh}
  \item \emph{Say it:} \emph{mimḥāh}, once more
  \item \emph{Say it:} say \emph{mimḥāh}
  \item \emph{Recall:} say the Arabic for a factory, then the Arabic for a square, then say \emph{mimḥāh} again
\end{itemize}

\begin{tcolorbox}[breakable,colback=teal!4,colframe=teal!35!black,title={Before you move on}]
What does \ar{ممحاة} mean? (\textquotedblleft{}An eraser\textquotedblright{}.) Say it once more.
\end{tcolorbox}
\section[\texorpdfstring{\ar{صف} (ṣaff)}{صف (ṣaff)}]{\ar{صف} (ṣaff) — a class}

\label{lesson:AR-C143-saff}



\begin{quote}
Before the new one: say the Arabic for a square, then the Arabic for an eraser.
\end{quote}

\subsection*{You'll want to know: \ar{صف}}
\textbf{\ar{صف}} — \emph{ṣaff} — \textquotedblleft{}a class\textquotedblright{}.

\begin{cousinweb}
\textbf{\ar{صف}} (\emph{ṣaff}) is built on the native Arabic root \emph{ṣ-f-f}, \textquotedblleft{}a row, a rank\textquotedblright{} (Lane). \emph{Ṣaff} is a row; the class is a row of pupils.
\end{cousinweb}

\begin{grammarlens}[title={What you've built}]
One more word for this chapter.
\end{grammarlens}

\subsection*{Your turn}
\begin{itemize}
\raggedright
  \item \emph{Say it:} \emph{ṣaff}
  \item \emph{Say it:} \emph{ṣaff}, once more
  \item \emph{Say it:} say \emph{ṣaff}
  \item \emph{Recall:} say the Arabic for a square, then the Arabic for an eraser, then say \emph{ṣaff} again
\end{itemize}

\begin{tcolorbox}[breakable,colback=teal!4,colframe=teal!35!black,title={Before you move on}]
What does \ar{صف} mean? (\textquotedblleft{}A class\textquotedblright{}.) Say it once more.
\end{tcolorbox}
\section[\texorpdfstring{\ar{طالب} (ṭālib)}{طالب (ṭālib)}]{\ar{طالب} (ṭālib) — a student}

\label{lesson:AR-C143-talib}



\begin{quote}
Before the new one: say the Arabic for an eraser, then the Arabic for a class.

Say \textquotedblleft{}summer\textquotedblright{} and \textquotedblleft{}winter.\textquotedblright{} (\emph{Ṣayf}; \emph{shitāʾ}.)
\end{quote}

\subsection*{You'll want to know: \ar{طالب}}
\textbf{\ar{طالب}} — \emph{ṭālib} — \textquotedblleft{}a student\textquotedblright{}.

\begin{cousinweb}
\textbf{\ar{طالب}} (\emph{ṭālib}) is built on the native Arabic root \emph{ṭ-l-b}, \textquotedblleft{}to seek, to ask for\textquotedblright{} (Lane). \emph{Ṭālib} is the active participle, \textquotedblleft{}one who seeks\textquotedblright{}.
\end{cousinweb}

\begin{grammarlens}[title={What you've built}]
One more word for this chapter.
\end{grammarlens}

\subsection*{Your turn}
\begin{itemize}
\raggedright
  \item \emph{Say it:} \emph{ṭālib}
  \item \emph{Say it:} \emph{ṭālib}, once more
  \item \emph{Say it:} say \emph{ṭālib}
  \item \emph{Recall:} say the Arabic for an eraser, then the Arabic for a class, then say \emph{ṭālib} again
\end{itemize}

\begin{tcolorbox}[breakable,colback=teal!4,colframe=teal!35!black,title={Before you move on}]
What does \ar{طالب} mean? (\textquotedblleft{}A student\textquotedblright{}.) Say it once more.
\end{tcolorbox}
\section[\texorpdfstring{\ar{وظيفة} (waẓīfa)}{وظيفة (waẓīfa)}]{\ar{وظيفة} (waẓīfa) — a job}

\label{lesson:AR-C143-wazifa}



\begin{quote}
Before the new one: say the Arabic for a class, then the Arabic for a student.
\end{quote}

\subsection*{You'll want to know: \ar{وظيفة}}
\textbf{\ar{وظيفة}} — \emph{waẓīfa} — \textquotedblleft{}a job\textquotedblright{}.

\begin{cousinweb}
\textbf{\ar{وظيفة}} (\emph{waẓīfa}) is built on the native Arabic root \emph{w-ẓ-f}, \textquotedblleft{}a daily allowance\textquotedblright{} (Lane). Lane's \emph{waẓīfa} is a daily allowance or fixed portion; the job is a later sense.
\end{cousinweb}

\begin{grammarlens}[title={What you've built}]
One more word for this chapter.
\end{grammarlens}

\subsection*{Your turn}
\begin{itemize}
\raggedright
  \item \emph{Say it:} \emph{waẓīfa}
  \item \emph{Say it:} \emph{waẓīfa}, once more
  \item \emph{Say it:} say \emph{waẓīfa}
  \item \emph{Recall:} say the Arabic for a class, then the Arabic for a student, then say \emph{waẓīfa} again
\end{itemize}

\begin{tcolorbox}[breakable,colback=teal!4,colframe=teal!35!black,title={Before you move on}]
What does \ar{وظيفة} mean? (\textquotedblleft{}A job\textquotedblright{}.) Say it once more.
\end{tcolorbox}
\section[\texorpdfstring{\ar{راتب} (rātib)}{راتب (rātib)}]{\ar{راتب} (rātib) — a salary}

\label{lesson:AR-C143-ratib}



\begin{quote}
Before the new one: say the Arabic for a student, then the Arabic for a job.
\end{quote}

\subsection*{You'll want to know: \ar{راتب}}
\textbf{\ar{راتب}} — \emph{rātib} — \textquotedblleft{}a salary\textquotedblright{}.

\begin{cousinweb}
\textbf{\ar{راتب}} (\emph{rātib}) is built on the native Arabic root \emph{r-t-b}, \textquotedblleft{}to be fixed, to be firm\textquotedblright{} (Lane). Lane's \emph{rātib} is a thing fixed and constant; the salary is a fixed sum.
\end{cousinweb}

\begin{grammarlens}[title={What you've built}]
One more word for this chapter.
\end{grammarlens}

\subsection*{Your turn}
\begin{itemize}
\raggedright
  \item \emph{Say it:} \emph{rātib}
  \item \emph{Say it:} \emph{rātib}, once more
  \item \emph{Say it:} say \emph{rātib}
  \item \emph{Recall:} say the Arabic for a student, then the Arabic for a job, then say \emph{rātib} again
\end{itemize}

\begin{tcolorbox}[breakable,colback=teal!4,colframe=teal!35!black,title={Before you move on}]
What does \ar{راتب} mean? (\textquotedblleft{}A salary\textquotedblright{}.) Say it once more.
\end{tcolorbox}
